\begin{abstract}
Linear-time recurrent sequence models such as Mamba-2 and Gated DeltaNet are replacing most attention
layers in recent large language models because they keep a fixed-size state instead of a key--value cache
that grows with the context. This report studies what that \emph{fading memory} costs and what it buys.
We implement, from scratch, a selective state-space model (SSM), Gated DeltaNet and a Transformer baseline,
each with a reference recurrence and a verified chunkwise-parallel training algorithm, plus two recent
extensions: a rotational (complex-valued, Mamba-3-style) SSM and DeltaProduct. All models are trained on a
single RTX~4050 laptop GPU on synthetic pattern-recognition probes: state tracking in the groups
$\mathbb{Z}_2$, $\mathbb{Z}_3$ and $S_3$ with length generalisation, and multi-query associative recall.
Three results stand out. (i) Allowing negative eigenvalues makes the SSM solve parity at 8$\times$ the
training length (0.999 vs.\ 0.546). (ii) DeltaProduct, a product of two generalised Householder
reflections per token, is the only model that learns the non-abelian group $S_3$ on all seeds and keeps
0.912 accuracy at 8$\times$ the training length, versus 0.542 for Gated DeltaNet. (iii) Recall degrades
exactly when the number of stored pairs outgrows the recurrent state, while Gated DeltaNet stores 64 pairs
with half the state of an SSM that fails. We also document failures: models that fit the training length
perfectly but learn shortcuts, a rotational SSM that can represent $\mathbb{Z}_3$ but does not learn it
(confirmed by reading its angles from the weights), and an under-budgeted recall sweep that ranked models
wrongly. A GPU study shows that, at feasible lengths, kernel quality dominates asymptotic complexity.
\end{abstract}
